\documentclass{amsart}
% For dealing with references we use the comment environment
\usepackage{verbatim}
\newenvironment{reference}{\comment}{\endcomment}
%\newenvironment{reference}{}{}
\newenvironment{slogan}{\comment}{\endcomment}
\newenvironment{history}{\comment}{\endcomment}

% For commutative diagrams we use Xy-pic
\usepackage[all]{xy}

% We use 2cell for 2-commutative diagrams.
\xyoption{2cell}
\UseAllTwocells

% We use multicol for the list of chapters between chapters
\usepackage{multicol}

% This is generally recommended for better output
\usepackage{lmodern}
\usepackage[T1]{fontenc}

% For cross-file-references

% Package for hypertext links:
\usepackage{hyperref}

% For any local file, say "hello.tex" you want to link to please

% Theorem environments.
%
\theoremstyle{plain}
\newtheorem{theorem}[subsection]{Theorem}
\newtheorem{proposition}[subsection]{Proposition}
\newtheorem{lemma}[subsection]{Lemma}

\theoremstyle{definition}
\newtheorem{definition}[subsection]{Definition}
\newtheorem{example}[subsection]{Example}
\newtheorem{exercise}[subsection]{Exercise}
\newtheorem{situation}[subsection]{Situation}

\theoremstyle{remark}
\newtheorem{remark}[subsection]{Remark}
\newtheorem{remarks}[subsection]{Remarks}

\numberwithin{equation}{subsection}

% Macros
%
\def\lim{\mathop{\mathrm{lim}}\nolimits}
\def\colim{\mathop{\mathrm{colim}}\nolimits}
\def\Spec{\mathop{\mathrm{Spec}}}
\def\Hom{\mathop{\mathrm{Hom}}\nolimits}
\def\Ext{\mathop{\mathrm{Ext}}\nolimits}
\def\SheafHom{\mathop{\mathcal{H}\!\mathit{om}}\nolimits}
\def\SheafExt{\mathop{\mathcal{E}\!\mathit{xt}}\nolimits}
\def\Sch{\mathit{Sch}}
\def\Mor{\mathop{\mathrm{Mor}}\nolimits}
\def\Ob{\mathop{\mathrm{Ob}}\nolimits}
\def\Sh{\mathop{\mathit{Sh}}\nolimits}
\def\NL{\mathop{N\!L}\nolimits}
\def\CH{\mathop{\mathrm{CH}}\nolimits}
\def\proetale{{pro\text{-}\acute{e}tale}}
\def\etale{{\acute{e}tale}}
\def\QCoh{\mathit{QCoh}}
\def\Ker{\mathop{\mathrm{Ker}}}
\def\Im{\mathop{\mathrm{Im}}}
\def\Coker{\mathop{\mathrm{Coker}}}
\def\Coim{\mathop{\mathrm{Coim}}}

% Boxtimes
%
\DeclareMathSymbol{\boxtimes}{\mathbin}{AMSa}{"02}

%
% Macros for moduli stacks/spaces
%
\def\QCohstack{\mathcal{QC}\!\mathit{oh}}
\def\Cohstack{\mathcal{C}\!\mathit{oh}}
\def\Spacesstack{\mathcal{S}\!\mathit{paces}}
\def\Quotfunctor{\mathrm{Quot}}
\def\Hilbfunctor{\mathrm{Hilb}}
\def\Curvesstack{\mathcal{C}\!\mathit{urves}}
\def\Polarizedstack{\mathcal{P}\!\mathit{olarized}}
\def\Complexesstack{\mathcal{C}\!\mathit{omplexes}}
% \Pic is the operator that assigns to X its picard group, usage \Pic(X)
% \Picardstack_{X/B} denotes the Picard stack of X over B
% \Picardfunctor_{X/B} denotes the Picard functor of X over B
\def\Pic{\mathop{\mathrm{Pic}}\nolimits}
\def\Picardstack{\mathcal{P}\!\mathit{ic}}
\def\Picardfunctor{\mathrm{Pic}}
\def\Deformationcategory{\mathcal{D}\!\mathit{ef}}

\setlength{\emergencystretch}{3em}
\begin{document}
\title[Stacks Project, Tag 02H6]{464 --- Smooth morphisms are exactly locally finitely presented formally smooth morphisms}
\maketitle
\noindent Copyright (C) 2005--2025 Johan de Jong. Stacks Project authors. Source: \href{https://stacks.math.columbia.edu/tag/02H6}{Tag 02H6}.

\noindent Editorial difficulty note (not author text). Difficulty H1: Affine formal lifting does not automatically glue to a scheme morphism on the whole affine test thickening. The proof constructs the sheaf of lifts as a torsor under the internal Hom of differentials and conormal modules, proves local solvability, and uses quasi-coherent H1 vanishing to obtain a global lift..

\noindent Editorial provenance note (not author text). History: excerpt from revision \texttt{a0444\allowbreak 6e57e\allowbreak c1fbc\allowbreak 252a8\allowbreak 71afc\allowbreak ec775\allowbreak 2fb28\allowbreak 07b14}, more-morphisms.tex, lines 3189--3282. Reference presentation was adapted; original-author context and core support blocks were retained or added. The mathematical wording is unchanged. Licensed under GNU FDL 1.2 or later; no invariant sections or cover texts. See ../licenses/COPYING. Context blocks reproduce author definitions and supporting results; they are not additional counted problems.

\begin{definition}
\label{cohomology-definition-torsor}
Let $X$ be a topological space.
Let $\mathcal{G}$ be a sheaf of (possibly non-commutative) groups on $X$.
A {\it pseudo torsor}, or more precisely a {\it pseudo $\mathcal{G}$-torsor},
is a sheaf of sets $\mathcal{F}$ on $X$ endowed with an action
$\mathcal{G} \times \mathcal{F} \to \mathcal{F}$ such that
\begin{enumerate}
\item whenever $\mathcal{F}(U)$ is nonempty the action
$\mathcal{G}(U) \times \mathcal{F}(U) \to \mathcal{F}(U)$
is simply transitive
\end{enumerate}
A {\it morphism of pseudo $\mathcal{G}$-torsors} $\mathcal{F} \to \mathcal{F}'$
is a morphism of sheaves of sets compatible with the
$\mathcal{G}$-actions. A {\it torsor}, or more precisely a
{\it $\mathcal{G}$-torsor}, is a pseudo $\mathcal{G}$-torsor such
that in addition
\begin{enumerate}
\item[(2)] for every $x \in X$ the stalk $\mathcal{F}_x$ is nonempty.
\end{enumerate}
A {\it morphism of $\mathcal{G}$-torsors} is a morphism of
pseudo $\mathcal{G}$-torsors. The {\it trivial $\mathcal{G}$-torsor}
is the sheaf $\mathcal{G}$ endowed with the obvious left
$\mathcal{G}$-action.
\end{definition}

\begin{definition}
\label{definition-formally-smooth}
Let $f : X \to S$ be a morphism of schemes.
We say $f$ is {\it formally smooth} if given any solid commutative diagram
$$
\xymatrix{
X \ar[d]_f & T \ar[d]^i \ar[l] \\
S & T' \ar[l] \ar@{-->}[lu]
}
$$
where $T \subset T'$ is a first order thickening of affine schemes over $S$
there exists a dotted arrow making the diagram commute.
\end{definition}

\begin{lemma}[Infinitesimal lifting criterion]
\label{lemma-smooth-formally-smooth}
Let $f : X \to S$ be a morphism of schemes.
The following are equivalent:
\begin{enumerate}
\item The morphism $f$ is smooth, and
\item the morphism $f$ is locally of finite presentation and
formally smooth.
\end{enumerate}
\end{lemma}

\begin{proof}
Assume $f : X \to S$ is locally of finite presentation and formally smooth.
Consider a pair of affine opens $\Spec(A) = U \subset X$ and
$\Spec(R) = V \subset S$
such that $f(U) \subset V$. By Lemma \href{https://stacks.math.columbia.edu/tag/02H3}{lemma formally smooth on opens (Tag 02H3)}
we see that $U \to V$ is formally smooth. By Lemma
\href{https://stacks.math.columbia.edu/tag/02H4}{lemma affine formally smooth (Tag 02H4)} we see that $R \to A$ is formally
smooth. By
Morphisms, Lemma \href{https://stacks.math.columbia.edu/tag/01TQ}{lemma locally finite presentation characterize (Tag 01TQ)}
we see that $R \to A$ is of finite presentation.
By Algebra, Proposition \href{https://stacks.math.columbia.edu/tag/00TN}{proposition smooth formally smooth (Tag 00TN)}
we see that $R \to A$ is smooth.
Hence by the definition of a smooth morphism we see that $X \to S$ is smooth.

\medskip\noindent
Conversely, assume that $f : X \to S$ is smooth. Consider a solid commutative
diagram
$$
\xymatrix{
X \ar[d]_f & T \ar[d]^i \ar[l]^a \\
S & T' \ar[l] \ar@{-->}[lu]
}
$$
as in Definition \ref{definition-formally-smooth}.
We will show the dotted arrow exists thereby
proving that $f$ is formally smooth.

\medskip\noindent
Let $\mathcal{F}$ be the sheaf of sets on $T'$ of Lemma \href{https://stacks.math.columbia.edu/tag/04FH}{lemma sheaf (Tag 04FH)}
in the special case discussed in Remark \href{https://stacks.math.columbia.edu/tag/04FK}{remark special case (Tag 04FK)}.
Let
$$
\mathcal{H} =
\SheafHom_{\mathcal{O}_T}(a^*\Omega_{X/S}, \mathcal{C}_{T/T'})
$$
be the sheaf of $\mathcal{O}_T$-modules with action
$\mathcal{H} \times \mathcal{F} \to \mathcal{F}$ as in
Lemma \href{https://stacks.math.columbia.edu/tag/04FJ}{lemma action sheaf (Tag 04FJ)}. Our goal is simply
to show that $\mathcal{F}(T) \not = \emptyset$. In other words we
are trying to show that $\mathcal{F}$ is a trivial $\mathcal{H}$-torsor
on $T$ (see Cohomology, Section \href{https://stacks.math.columbia.edu/tag/02FN}{section h1 torsors (Tag 02FN)}).
There are two steps: (I) To show that $\mathcal{F}$ is a torsor
we have to show that $\mathcal{F}_t \not = \emptyset$ for all $t \in T$ (see
Cohomology, Definition \ref{cohomology-definition-torsor}).
(II) To show that $\mathcal{F}$ is the trivial torsor it suffices
to show that $H^1(T, \mathcal{H}) = 0$ (see
Cohomology, Lemma \href{https://stacks.math.columbia.edu/tag/02FQ}{lemma torsors h1 (Tag 02FQ)} --
we may use either cohomology
of $\mathcal{H}$ as an abelian sheaf or as an $\mathcal{O}_T$-module,
see Cohomology, Lemma \href{https://stacks.math.columbia.edu/tag/01F1}{lemma modules abelian (Tag 01F1)}).

\medskip\noindent
First we prove (I). To see this, for every $t \in T$ we can
choose an affine open $U \subset T$ neighbourhood of $t$
such that $a(U)$ is contained
in an affine open $\Spec(A) = W \subset X$
which maps to an affine open $\Spec(R) = V \subset S$.
By Morphisms, Lemma \href{https://stacks.math.columbia.edu/tag/01V6}{lemma smooth characterize (Tag 01V6)}
the ring map $R \to A$ is smooth.
Hence by Algebra, Proposition \href{https://stacks.math.columbia.edu/tag/00TN}{proposition smooth formally smooth (Tag 00TN)}
the ring map $R \to A$ is formally smooth.
Lemma \href{https://stacks.math.columbia.edu/tag/02H4}{lemma affine formally smooth (Tag 02H4)}
in turn implies that $W \to V$ is formally smooth.
Hence we can lift $a|_U : U \to W$ to a $V$-morphism
$a' : U' \to W \subset X$ showing that $\mathcal{F}(U) \not = \emptyset$.

\medskip\noindent
Finally we prove (II).
By Morphisms, Lemma \href{https://stacks.math.columbia.edu/tag/01V3}{lemma finite presentation differentials (Tag 01V3)}
we see that $\Omega_{X/S}$ is of finite presentation
(it is even finite locally free by
Morphisms, Lemma \href{https://stacks.math.columbia.edu/tag/02G1}{lemma smooth omega finite locally free (Tag 02G1)}).
Hence $a^*\Omega_{X/S}$ is of finite presentation (see
Modules, Lemma \href{https://stacks.math.columbia.edu/tag/01BQ}{lemma pullback finite presentation (Tag 01BQ)}).
Hence the sheaf
$\mathcal{H} =
\SheafHom_{\mathcal{O}_T}(a^*\Omega_{X/S}, \mathcal{C}_{T/T'})$
is quasi-coherent by the discussion in
Schemes, Section \href{https://stacks.math.columbia.edu/tag/01LA}{section quasi coherent (Tag 01LA)}.
Thus by Cohomology of Schemes, Lemma
\href{https://stacks.math.columbia.edu/tag/01XB}{lemma quasi coherent affine cohomology zero (Tag 01XB)}
we have $H^1(T, \mathcal{H}) = 0$ as desired.
\end{proof}
\end{document}
